%=====================================================================
%  Front matter: abstract, contents, lists, abbreviations
%=====================================================================
\chapter*{Abstract}
\addcontentsline{toc}{chapter}{Abstract}

Equity markets are driven simultaneously by structured price and volume records, by
qualitative information carried in news and social media, and by slower macroeconomic
conditions. Each stream is individually incomplete, and each becomes informative or
uninformative at different times. Conventional forecasting systems compress this
heterogeneity into a single numerical pipeline, report a statistical accuracy figure, and
stop short of the decision the investor actually has to take. This research addresses that
shortfall by developing, implementing and validating an integrated framework in which
multi-source fusion, language-model-based sentiment extraction, transparent reasoning and
decision-level allocation are treated as one connected engineering problem rather than as
four separate ones.

The first objective develops a comprehensive model integrating financial metrics, news
sentiment, social media sentiment and macroeconomic indicators. A hierarchical
Multi-Source Fusion Network is designed in which convolutional encoders process indicator
streams, transformer encoders process text, engineered temporal features carry
macroeconomic conditions, and a cross-modal attention layer assigns time-varying weight to
each source before a bidirectional recurrent encoder produces the forecast. Evaluated on
ten years of National Stock Exchange data across ten sectoral indices, the fused model
attains a normalised mean absolute error of 0.018, a coefficient of determination of 0.988
and a directional accuracy of 84.7 per cent at the index level, against 76.3 per cent for
a simple averaging ensemble over identical inputs. Removing both textual sources costs
10.9 per cent of explained variance, exceeding the sum of their individual removals and
confirming that sentiment carries information that price history does not. The fusion
mechanism is then strengthened by a Bayesian formulation in which the weight assigned to a
source is derived from its own recent predictive uncertainty through a softmax over
negative error variances, so that a stream is discounted at the moment it becomes
unreliable rather than at the end of a training cycle. On the Nifty~50 this dynamic rule
improves the Sharpe ratio by 22.7 per cent and reduces maximum drawdown by 18.3 per cent
against the strongest static-fusion baseline, while the same expert networks fused with
fixed weights perform no better than that baseline, which locates the gain in the weighting
mechanism rather than in the networks.

The second objective designs a hybrid framework coupling large language models with
conventional financial indicators. A transformer sentiment encoder converts news headlines
into a continuous polarity score, which is time-aligned to trading sessions and fused with
fifteen technical indicators before gradient-boosted classification. On Indian equity data
the sentiment score emerges as the highest-gain predictor, correlating 0.74 with momentum
yet only 0.04 with price, and the sentiment-augmented model raises directional accuracy
from a 46.43 per cent baseline average to 47.72 per cent while producing the highest
prediction confidence of all models examined. The absolute level is reported without
embellishment: at single-name daily horizons all models cluster near the random
classifier, so language-model sentiment supplies a measurable but bounded increment rather
than a decisive edge.

The third objective establishes transparency and examines sentiment noise and
sentiment-to-price temporal dynamics. A Dynamic-Causal Hybrid Framework replaces
correlation-based feature selection with a directed causal adjacency matrix recovered by
conditional independence testing, encodes intraday, daily and weekly horizons through a
multi-scale memory encoder, and allocates capital through a credibilistic optimiser. On
308 trading sessions of NIFTY50 data it recovers a causal graph of density 0.850 with an
identifiable causal hub, and demonstrates that strong causal influence can coexist with
weak linear correlation --- the mechanistic justification for preferring causal to
correlational feature selection. Sentiment quality is then treated as a measurable
property rather than an assumption: a sentiment composite is validated against four
independently published series, carries the expected sign against every one of them, and
separates the identified market states at $p = 1.8 \times 10^{-8}$.

The fourth objective evaluates and validates the framework as a decision system. Latent
market states are identified by an expanding-window Gaussian mixture model, the fused
signal enters a constrained quadratic program with covariance shrinkage and a turnover
penalty, and the resulting portfolio is scaled by a regime-dependent risk budget. Over
eighteen years of ten liquid multi-asset instruments under a walk-forward protocol with
transaction costs, the framework attains a Sharpe ratio of 0.878 against 0.628 for an
equal-weighted benchmark and reduces maximum drawdown from $-36.1$ to $-19.0$ per cent. It
outperforms in all four stress episodes, wins on drawdown in eighteen of nineteen calendar
years and in all eighteen specifications examined, and on a 665-session holdout that
postdates every design decision it records a higher Sharpe ratio and a smaller drawdown
than all three benchmarks. Ablation at matched average exposure isolates the source of the
gain: timing risk by market state improves both the Sharpe ratio and the drawdown with no
change in mean return, whereas the signal fusion layer contributes nothing detectable to
cross-sectional selection. Cross-market transfer to a concentrated mega-capitalisation
universe reproduces the defensive behaviour but loses on risk-adjusted return, which bounds
applicability to universes with genuine dispersion in asset volatility.

Taken together, the six publications underlying this report demonstrate a coherent
progression from information integration, through language-model-based signal extraction
and transparent reasoning, to validated decision making. The component-level nulls are
retained rather than removed, because they identify precisely which layers of a fusion
architecture earn their place and which do not.

\vspace{0.4em}
\noindent\textbf{Keywords:} multi-source data fusion; sentiment analysis; large language
models; explainable artificial intelligence; regime detection; portfolio construction;
Indian equity market.

%---------------------------------------------------------------------
\clearpage
\renewcommand{\contentsname}{Table of Contents}
\tableofcontents

\clearpage
\renewcommand{\listfigurename}{List of Figures}
\listoffigures

\clearpage
\renewcommand{\listtablename}{List of Tables}
\listoftables

%---------------------------------------------------------------------
\clearpage
\chapter*{Abbreviations}
\addcontentsline{toc}{chapter}{Abbreviations}

\begin{center}
\footnotesize
\setstretch{1.0}
\begin{tabular}{@{}p{1.9cm}p{4.5cm}p{1.9cm}p{4.9cm}@{}}
ANOVA & Analysis of Variance & MAE & Mean Absolute Error\\
ARIMA & Autoregressive Integrated Moving Average & MDD & Maximum Drawdown\\
AUC & Area Under the Curve & ML & Machine Learning\\
Bi-LSTM & Bidirectional Long Short-Term Memory & MLP & Multi-Layer Perceptron\\
BSE & Bombay Stock Exchange & MOSPI & Ministry of Statistics and Programme Implementation\\
CAGR & Compound Annual Growth Rate & MSFN & Multi-Source Fusion Network\\
CIT & Conditional Independence Test & NLP & Natural Language Processing\\
CNN & Convolutional Neural Network & NSE & National Stock Exchange\\
CPI & Consumer Price Index & OHLCV & Open, High, Low, Close, Volume\\
CVaR & Conditional Value at Risk & QP & Quadratic Program\\
DA & Directional Accuracy & RBI & Reserve Bank of India\\
DCHF & Dynamic-Causal Hybrid Framework & RDM-DEA & Range Directional Measure Data Envelopment Analysis\\
DL & Deep Learning & RMSE & Root Mean Square Error\\
ETF & Exchange-Traded Fund & RSI & Relative Strength Index\\
GMM & Gaussian Mixture Model & SEBI & Securities and Exchange Board of India\\
GRU & Gated Recurrent Unit & SHAP & Shapley Additive Explanations\\
IC & Information Coefficient & SMA & Simple Moving Average\\
IIP & Index of Industrial Production & SVM & Support Vector Machine\\
LLM & Large Language Model & VADER & Valence Aware Dictionary and sEntiment Reasoner\\
LSTM & Long Short-Term Memory & VIX & Volatility Index\\
MACD & Moving Average Convergence Divergence & WPI & Wholesale Price Index\\
& & XAI & Explainable Artificial Intelligence\\
\end{tabular}
\end{center}
